\chapter{Introduction}
\label{ch:into} % This how you label a chapter and the key (e.g., ch:into) will be used to refer this chapter ``Introduction'' later in the report. 
% the key ``ch:into'' can be used with command \ref{ch:intor} to refere this Chapter.

%%%%%%%%%%%%%%%%%%%%%%%%%%%%%%%%%%%%%%%%%%%%%%%%%%%%%%%%%%%%%%%%%%%%%%%%%%%%%%%%%%%
\textit{Recommender Systems (RS)} are a large family of algorithms that are designed to narrow down the choice of potential options and surface the most relevant items to users. In recent decades globalisation has enabled retailers to offer a vast range of products to their customers. It has been previously shown that abundance of options, particularly when choice complexity is high (e.g., number of distinguishing attributes is large), may lead to user frustration and poor commercial outcomes \citep{greifeneder2010less}. Traditionally, this problem has been mitigated by human advisors, however, the widely recognised ‘death of the high street’ trend calls for the need to offer recommendation services online, in a commercially viable, scalable and automated way. In the last two decades this need has been addressed by RS, a family of algorithms designed to personalise user experience and increasingly adopted by e-commerce businesses \citep{chandrashekara2018fuzzy, kucherbaev2017chatbots, zhao2013interactive}. RS surface the most relevant options to users based on a range of data sources, which include behavioural (collaborative) data, user and product attributes. However, traditional recommender systems cannot replicate the experience of talking to a human advisor, as they lack interactivity, which leads to such problems as \textit{cold start}, where existing data is not sufficient to compute a recommendation, and \textit{taste drift}, where user’s requirements change over time, and recommendations computed based on historical data become irrelevant. 

Intelligent conversational systems have high potential in commercial settings. A combination of \textit{RS} with intelligent conversational algorithms has resulted in the emergence of the \textit{Conversational Recommender Systems (CRS)}, a subset of RS that are interactive in their nature. \textit{CRS} build upon traditional offline algorithms by introducing interactivity, which results in a capability for active \textit{human-in-the-loop} learning.  \textit{CRS} provide users with an opportunity to communicate their preferences and utilise this information to update user representations in real time.  \textit{CRS} can identify gaps in user representations and proactively request information from users to compute more appropriate recommendations. This closely mimics a conversation between two humans, particularly where natural language is used as an interface.

CRS is a relatively new field, and very few of such algorithms have been used or tested in real world settings. Laboratory prototypes range from adaptive personalised web interfaces \citep{ayundhita2019ontology} to natural language based chat bots \citep{sun2018conversational, li2018towards} to voice based systems \citep{dalton2018vote}. Despite the differences in the mode of interaction, the main goal of all such systems is to elicit information from users in order to optimise recommendations in a way that increases their relevance. Existing CRS algorithms that are optimised for recommendation relevance are all based on synthetic datasets, which results in contrived and scripted conversations that do not generalise well to real-world scenarios \citep{jannach2021survey}. In contrast, works that focus on training agents capable of sustaining natural dialogues do not convincingly show that the resulting CRS produce quality recommendation \citep{li2018towards, kang2019recommendation}. 

Recent survey of the state-of-the-art of the CRS field by \cite{jannach2022conversational} revealed that most natural language based systems that have practical application are based on predefined templates and rule-based logic. Laboratory prototypes trained end-to-end to natural datasets perform worse in human evaluation tests and have a high rate of unusable responses, where model fails to correctly identify and address user intent.

In addition to that, maintaining a multi-turn and mixed-initiative conversations remains a challenge. \cite{jannach2022conversational} note that in order to mimic a dialogue between two humans, CRS must support user-driven, agent-driven and mixed initiative interactions, as well as respond to a range of intents including providing or revising preference information, asking for additional information, rejecting a recommendation, including intents not directly related to recommendation problem, such as greetings or phatic expressions. Authors report that majority of algorithms reported in the literature focus only on a small subset of these intents, which limits their usabilty in practical terms.

Evaluation of CRS, particularly those with natural language interface remains a challenge as well \cite{jannach2022conversational, jannach2022evaluating}. CRS tend to complex multi-component systems which complicates their direct comparison. The recommendation component traditionally is evaluated using metrics related to item prediction (MRR, AUC, Precision@$k$ etc.), whereas true recommendation problem is more complex and involve aspects such as serendipity and surprise. Automated metrics related to evaluating quality of generated utterances are also flawed, as metrics such as BLEU have poor correlation to human evaluation. Some studies evaluate metrics such as success rate and rejection rate based on user simulations, however simulated user may not behave realistically. Human evaluation of metrics such as decision confidence, perceived recommendation quality, purchase intentions and user satisfaction remain the most informative, despite being highly subjective.

Building upon prior research in the field of CRS, this thesis aims to construct a comprehensive and adaptable algorithm for CRS with an emphasis on devising an efficient strategy for extracting valuable information from users within the vast and diverse realm of potential item features. This thesis focuses on addressing a significant issue in the field of dialogue systems: the absence of natural language interfaces that can effectively elicit information from users and generate accurate recommendations. Currently available systems exhibit either a natural sounding dialogue but lack satisfactory recommendation capabilities, or they prioritise generation of high-quality recommendations at the expense of using rigid templates that result in a synthetic conversation. The proposed algorithm design will effectively utilise multi-modal recommendation data, encompassing collaborative behavioral data (such as user-item rating matrices) and item attributes. By doing so, the system will be capable of effectively eliciting user preferences in a human-like manner and leveraging this information to generate high-quality recommendations. The algorithm will harness its interactive capabilities to optimise recommendation relevance, aid users in making satisfactory consumption decisions, and have a potential to enhance conversion rates for businesses. 

To address the aforementioned research challenges, the following thesis objectives have been established:

\begin{itemize}
    \item  To design a recommender system specifically tailored for dialogue scenarios where there is limited user information. The goal is to develop a system that can effectively generate accurate recommendations despite the limited initial user input.
    \item  To leverage the interactivity of the CRS to proactively extract crucial information from users. By actively engaging with the users, the system aims to efficiently gather the most relevant information necessary for making optimal recommendations in a timely manner.
    \item  To create a flexible and human-like interface for the recommender system. The objective is to enhance the user experience by designing an interface that mimics natural human conversation, allowing for a more intuitive and engaging interaction between the user and the system.
\end{itemize}

To achieve these goals, this thesis explored existing state-of-the-art methods in CRS space in order to to study the proposed system architectures, methodologies, as well as analyse strengths and weaknesses of different approaches. In addition to that, it explores and examines state-of-the-art methods in the adjacent fields of task-oriented conversational systems, recommender systems, as well as natural language understanding and generation. A Recommendation module is designed especially for the dialogue scenario and trained based on a collaborative dataset augmented it with item attribute data. An Information Elicitation module is designed and trained through bot-play, enabling it to learn from data without relying on synthetic training environments. The algorithm performance is evaluated using experiments based on simulated user profiles.

The significance of the proposed work is primarily applied. The thesis will lean on the state-of-the-art methodology developed in the domains of Reinforcement Learning and Recommender Systems in order to address the unique challenges involved in designing a CRS, a task-oriented system that aims to emulate the experience of conversing with a human agent. 

From an academic standpoint, this thesis aims to add to the growing body of research in the emerging field of CRS. Specifically, it will seek to contribute to information elicitation methods by applying bot-play training protocol to multi-modal and high-dimensional space of recommendation data. 

From an industry standpoint, the thesis aims to narrow down the gap between artificial CRS and human advisors. An automated and scalable system crucial in order for e-commerce retailers to offer assistance to users, given the vast range of products on offer, which can reach millions of product titles, the scale of the operation, which caters to hundreds of millions of users, and low profit margins. While traditional RS partially address this problem, their lack of interactivity does not allow to directly query users for information that is necessary to generate a useful recommendation. Although there have been numerous attempts in recent years to design CRS, very few of them have been tested or applied outside of laboratory \citep{mahmood2009improving}, and none of those utilise natural language interface. The ones that do \citep{li2018towards, kang2019recommendation}, typically focus on the quality of the generated utterances and not on recommendation relevance, which remains the main driver of customer conversion. This thesis aims to design an end-to-end system applicable to different industries and capable of maintaining a conversation with users in a way that helps them share their requirements and preferences, and process this information in order to narrow down the range of choices and recommend products that are highly relevant to users, which, in turn, is expected to result in increased conversion rates for the businesses. 

The main contributions of this thesis is a practically applicable CRS algorithm that converses with users, builds accurate representations of their preferences and requirements and recommends satisfactory item choices. Further to that, this thesis describes the first attempt to apply bot-play as choice of a training approach in order to optimise the quality of recommendations.

The remaining sections of this thesis are structured into five chapters. Chapter 2 provides an extensive review of the existing literature on Conversational Recommender Systems, encompassing recent advancements and research in this domain. Chapter 3 focuses on the exploration of reinforcement learning techniques utilised for training such systems, examining their efficacy and applicability. In Chapter 4, the recommendation module is meticulously examined, introducing a novel architectural design that incorporates a Long Short-Term Memory (LSTM) network with attention mechanism. This innovative approach adopts an item-to-item approach to predict movie ratings. Chapter 5 proposes a reinforcement learning technique known as 'bot-play' to improve the elicitation of user preferences. This technique involves the proactive interaction between a QBot and an ABot, with the QBot prompting the ABot to share its preferences and modelling the user profile. Finally, Chapter 5 encompasses a comprehensive discussion of the conclusions drawn from the study and outlines potential avenues for future research and development.
